% GENERATED FILE. Edit canonical lessons, then run npm run generate:books.
% canonical-source-hash: fnv1a64:e76e9db65374046a
% canonical-lessons: TE-C381-tegincu, TE-C381-apaharincu, TE-C381-vicarincu, TE-C381-japtu-ceyu, TE-C381-dandayatra, TE-R381-review-374-381, TE-R381-review-final-381

\chapter{Daring Deeds}
\label{ch:te-daring-deeds}
\hlchaptermodality{381}

\begin{chapteropening}
\textbf{By the end of this chapter:} \emph{I can say to dare, to abduct, to interrogate, to confiscate and an invasion.}

Complete the last lesson of chapter 381: I can say to dare, to abduct, to interrogate, to confiscate and an invasion.
\end{chapteropening}

\section[\texorpdfstring{\te{తెగించు} (tegiñcu)}{తెగించు (tegiñcu)}]{\te{తెగించు} (tegiñcu) — to dare; to risk everything}

\label{lesson:TE-C381-tegincu}



\begin{quote}
Before the new one: say the Telugu for a dead body, then the Telugu for to haunt.
\end{quote}

\subsection*{You'll want to know: \te{తెగించు}}
\textbf{\te{తెగించు}} — \emph{tegiñcu} — \textquotedblleft{}to dare; to risk everything\textquotedblright{}.

\begin{quote}
\te{ఆ} \te{పిల్లవాడు} \te{తెగించి} \te{నదిలోకి} \te{దూకాడు}.
\end{quote}

\begin{quote}
\emph{ā pillavāḍu tegiñci nadilōki dūkāḍu.}
\end{quote}

\begin{quote}
The boy dared to jump into the river.
\end{quote}

Past (I …): \te{తెగించాను}, \emph{tegiñcānu}.

Often as \te{తెగించి}, daring: \te{తెగించి} \te{ముందుకు} \te{వెళ్ళు}.

\begin{cousinweb}
\textbf{\te{తెగించు}} (\emph{tegiñcu}) is native Dravidian: DEDR 3416 lists Telugu \emph{tegu}, \emph{tegincu}, \textquotedblleft{}to dare, venture\textquotedblright{}, beside \emph{tempu}, \textquotedblleft{}boldness\textquotedblright{}.
\end{cousinweb}

\begin{grammarlens}[title={What you've built}]
One more word for telling what happened.
\end{grammarlens}

\subsection*{Your turn}
\begin{itemize}
\raggedright
  \item \emph{Say it:} \emph{tegiñcu}
  \item \emph{Say it:} \emph{tegiñcu}, once more
  \item \emph{Say it:} \emph{ā pillavāḍu tegiñci nadilōki dūkāḍu.} — The boy dared to jump into the river.
  \item \emph{Recall:} say the Telugu for a dead body, then the Telugu for to haunt, then say \emph{tegiñcu} again
\end{itemize}

\begin{tcolorbox}[breakable,colback=teal!4,colframe=teal!35!black,title={Before you move on}]
What does \te{తెగించు} mean? (\textquotedblleft{}To dare; to risk everything\textquotedblright{}.) Say it once more.
\end{tcolorbox}
\section[\texorpdfstring{\te{అపహరించు} (apahariñcu)}{అపహరించు (apahariñcu)}]{\te{అపహరించు} (apahariñcu) — to abduct; to carry off}

\label{lesson:TE-C381-apaharincu}



\begin{quote}
Before the new one: say the Telugu for to haunt, then the Telugu for to dare.
\end{quote}

\subsection*{You'll want to know: \te{అపహరించు}}
\textbf{\te{అపహరించు}} — \emph{apahariñcu} — \textquotedblleft{}to abduct; to carry off\textquotedblright{}.

\begin{quote}
\te{దొంగలు} \te{ఒక} \te{వ్యాపారిని} \te{అపహరించారు}.
\end{quote}

\begin{quote}
\emph{doṅgalu oka vyāpārini apahariñcāru.}
\end{quote}

\begin{quote}
The robbers abducted a merchant.
\end{quote}

Past (I …): \te{అపహరించాను}, \emph{apahariñcānu}.

A formal word: to kidnap a person, or to steal goods.

\begin{cousinweb}
\textbf{\te{అపహరించు}} (\emph{apahariñcu}) is built on Sanskrit \emph{apaharaṇa}, \textquotedblleft{}taking away, carrying off; stealing\textquotedblright{} (Monier-Williams), with the Telugu verb ending \emph{-iñcu}.
\end{cousinweb}

\begin{grammarlens}[title={What you've built}]
One more word for telling what happened.
\end{grammarlens}

\subsection*{Your turn}
\begin{itemize}
\raggedright
  \item \emph{Say it:} \emph{apahariñcu}
  \item \emph{Say it:} \emph{apahariñcu}, once more
  \item \emph{Say it:} \emph{doṅgalu oka vyāpārini apahariñcāru.} — The robbers abducted a merchant.
  \item \emph{Recall:} say the Telugu for to haunt, then the Telugu for to dare, then say \emph{apahariñcu} again
\end{itemize}

\begin{tcolorbox}[breakable,colback=teal!4,colframe=teal!35!black,title={Before you move on}]
What does \te{అపహరించు} mean? (\textquotedblleft{}To abduct; to carry off\textquotedblright{}.) Say it once more.
\end{tcolorbox}
\section[\texorpdfstring{\te{విచారించు} (vicāriñcu)}{విచారించు (vicāriñcu)}]{\te{విచారించు} (vicāriñcu) — to interrogate; to inquire}

\label{lesson:TE-C381-vicarincu}



\begin{quote}
Before the new one: say the Telugu for to dare, then the Telugu for to abduct.
\end{quote}

\subsection*{You'll want to know: \te{విచారించు}}
\textbf{\te{విచారించు}} — \emph{vicāriñcu} — \textquotedblleft{}to interrogate; to inquire\textquotedblright{}.

\begin{quote}
\te{పోలీసులు} \te{నిందితుడిని} \te{గంటల} \te{తరబడి} \te{విచారించారు}.
\end{quote}

\begin{quote}
\emph{pōlīsulu nindituḍini gaṇṭala tarabaḍi vicāriñcāru.}
\end{quote}

\begin{quote}
The police questioned the accused for hours.
\end{quote}

Past (I …): \te{విచారించాను}, \emph{vicāriñcānu}.

You know \te{విచారం}, sorrow; the verb means to inquire, to look into.

\begin{cousinweb}
\textbf{\te{విచారించు}} (\emph{vicāriñcu}) is built on Sanskrit \emph{vicāra}, \textquotedblleft{}mode of action; pondering\textquotedblright{} (CDIAL 11683), with the Telugu verb ending \emph{-iñcu}.
\end{cousinweb}

\begin{grammarlens}[title={What you've built}]
One more word for telling what happened.
\end{grammarlens}

\subsection*{Your turn}
\begin{itemize}
\raggedright
  \item \emph{Say it:} \emph{vicāriñcu}
  \item \emph{Say it:} \emph{vicāriñcu}, once more
  \item \emph{Say it:} \emph{pōlīsulu nindituḍini gaṇṭala tarabaḍi vicāriñcāru.} — The police questioned the accused for hours.
  \item \emph{Recall:} say the Telugu for to dare, then the Telugu for to abduct, then say \emph{vicāriñcu} again
\end{itemize}

\begin{tcolorbox}[breakable,colback=teal!4,colframe=teal!35!black,title={Before you move on}]
What does \te{విచారించు} mean? (\textquotedblleft{}To interrogate; to inquire\textquotedblright{}.) Say it once more.
\end{tcolorbox}
\section[\texorpdfstring{\te{జప్తు} \te{చేయు} (japtu cēyu)}{జప్తు చేయు (japtu cēyu)}]{\te{జప్తు} \te{చేయు} (japtu cēyu) — to confiscate; to seize}

\label{lesson:TE-C381-japtu-ceyu}



\begin{quote}
Before the new one: say the Telugu for to abduct, then the Telugu for to interrogate.
\end{quote}

\subsection*{You'll want to know: \te{జప్తు} \te{చేయు}}
\textbf{\te{జప్తు} \te{చేయు}} — \emph{japtu cēyu} — \textquotedblleft{}to confiscate; to seize\textquotedblright{}.

\begin{quote}
\te{పోలీసులు} \te{దొంగ} \te{సొమ్మును} \te{జప్తు} \te{చేశారు}.
\end{quote}

\begin{quote}
\emph{pōlīsulu doṅga sommunu japtu cēśāru.}
\end{quote}

\begin{quote}
The police seized the stolen money.
\end{quote}

Past (I …): \te{జప్తు} \te{చేశాను}, \emph{japtu cēśānu}.

A seizure by law: officials \te{జప్తు} \te{చేస్తారు}.

\begin{cousinweb}
\textbf{\te{జప్తు} \te{చేయు}} (\emph{japtu cēyu}) is a loan from Persian \emph{zabt}, \textquotedblleft{}seizure\textquotedblright{}, with the native verb \emph{cēyu}, \textquotedblleft{}to do\textquotedblright{}.
\end{cousinweb}

\begin{grammarlens}[title={What you've built}]
One more word for telling what happened.
\end{grammarlens}

\subsection*{Your turn}
\begin{itemize}
\raggedright
  \item \emph{Say it:} \emph{japtu cēyu}
  \item \emph{Say it:} \emph{japtu cēyu}, once more
  \item \emph{Say it:} \emph{pōlīsulu doṅga sommunu japtu cēśāru.} — The police seized the stolen money.
  \item \emph{Recall:} say the Telugu for to abduct, then the Telugu for to interrogate, then say \emph{japtu cēyu} again
\end{itemize}

\begin{tcolorbox}[breakable,colback=teal!4,colframe=teal!35!black,title={Before you move on}]
What does \te{జప్తు} \te{చేయు} mean? (\textquotedblleft{}To confiscate; to seize\textquotedblright{}.) Say it once more.
\end{tcolorbox}
\section[\texorpdfstring{\te{దండయాత్ర} (daṇḍayātra)}{దండయాత్ర (daṇḍayātra)}]{\te{దండయాత్ర} (daṇḍayātra) — an invasion; a military campaign}

\label{lesson:TE-C381-dandayatra}



\begin{quote}
Before the new one: say the Telugu for to interrogate, then the Telugu for to confiscate.
\end{quote}

\subsection*{You'll want to know: \te{దండయాత్ర}}
\textbf{\te{దండయాత్ర}} — \emph{daṇḍayātra} — \textquotedblleft{}an invasion; a military campaign\textquotedblright{}.

\begin{quote}
\te{ఆ} \te{చక్రవర్తి} \te{చాలా} \te{దండయాత్రలు} \te{చేశాడు}.
\end{quote}

\begin{quote}
\emph{ā cakravarti cālā daṇḍayātralu cēśāḍu.}
\end{quote}

\begin{quote}
That emperor led many military campaigns.
\end{quote}

\te{దండయాత్ర} \te{చేయు}, to invade, to lead a campaign.

\begin{cousinweb}
\textbf{\te{దండయాత్ర}} (\emph{daṇḍayātra}) is borrowed from Sanskrit \emph{daṇḍayātrā}, \textquotedblleft{}a procession; a military expedition\textquotedblright{} (Monier-Williams), from \emph{daṇḍa}, \textquotedblleft{}a stick, a rod\textquotedblright{} (CDIAL 6128), and \emph{yātrā}, \textquotedblleft{}a journey\textquotedblright{} (CDIAL 10456).
\end{cousinweb}

\begin{grammarlens}[title={What you've built}]
One more word for telling what happened.
\end{grammarlens}

\subsection*{Your turn}
\begin{itemize}
\raggedright
  \item \emph{Say it:} \emph{daṇḍayātra}
  \item \emph{Say it:} \emph{daṇḍayātra}, once more
  \item \emph{Say it:} \emph{ā cakravarti cālā daṇḍayātralu cēśāḍu.} — That emperor led many military campaigns.
  \item \emph{Recall:} say the Telugu for to interrogate, then the Telugu for to confiscate, then say \emph{daṇḍayātra} again
\end{itemize}

\begin{tcolorbox}[breakable,colback=teal!4,colframe=teal!35!black,title={Before you move on}]
What does \te{దండయాత్ర} mean? (\textquotedblleft{}An invasion; a military campaign\textquotedblright{}.) Say it once more.
\end{tcolorbox}
\section[(dialogue)]{Rebels, legends and the press}

\label{lesson:TE-R381-review-374-381}



\begin{quote}
Say the Telugu for a commentary and a commentator.
\end{quote}

\subsection*{Your turn}
Say these aloud:

\begin{itemize}
\raggedright
  \item a commentary, a commentator, a listener, to behold, a face-to-face interview
  \item a journalist, printing, a sensation, to submit a report, to respond
  \item a response, a start, an ending, an episode, by chance
  \item a disaster, a commotion, protection, a patrol, an action
  \item a brawl, gunfire, to injure, to grapple, to gather round
  \item to throng, to fall down, a victim, to petition, to beg
  \item ago, before that, from then on, a dead body, to haunt
  \item to dare, to abduct, to interrogate, to confiscate, an invasion
\end{itemize}

\begin{tcolorbox}[breakable,colback=teal!4,colframe=teal!35!black,title={Before you move on}]
A commentary? (\textbf{\te{వ్యాఖ్యానం}}, \emph{vyākhyānaṁ}.) A commentator? (\textbf{\te{వ్యాఖ్యాత}}, \emph{vyākhyāta}.) A listener? (\textbf{\te{శ్రోత}}, \emph{śrōta}.) To behold? (\textbf{\te{తిలకించు}}, \emph{tilakiñcu}.) A face-to-face interview? (\textbf{\te{ముఖాముఖి}}, \emph{mukhāmukhi}.) A journalist? (\textbf{\te{పాత్రికేయుడు}}, \emph{pātrikēyuḍu}.) Printing? (\textbf{\te{ముద్రణ}}, \emph{mudraṇa}.) A sensation? (\textbf{\te{సంచలనం}}, \emph{sañcalanaṁ}.) To submit a report? (\textbf{\te{నివేదించు}}, \emph{nivēdiñcu}.) To respond? (\textbf{\te{స్పందించు}}, \emph{spandiñcu}.) A response? (\textbf{\te{స్పందన}}, \emph{spandana}.) A start? (\textbf{\te{ప్రారంభం}}, \emph{prārambhaṁ}.) An ending? (\textbf{\te{ముగింపు}}, \emph{mugimpu}.) An episode? (\textbf{\te{ఘట్టం}}, \emph{ghaṭṭaṁ}.) By chance? (\textbf{\te{యాదృచ్ఛికంగా}}, \emph{yādṛcchikaṅgā}.) A disaster? (\textbf{\te{విపత్తు}}, \emph{vipattu}.) A commotion? (\textbf{\te{అలజడి}}, \emph{alajaḍi}.) Protection? (\textbf{\te{రక్షణ}}, \emph{rakṣaṇa}.) A patrol? (\textbf{\te{గస్తీ}}, \emph{gastī}.) An action? (\textbf{\te{చర్య}}, \emph{carya}.) A brawl? (\textbf{\te{కొట్లాట}}, \emph{koṭlāṭa}.) Gunfire? (\textbf{\te{కాల్పులు}}, \emph{kālpulu}.) To injure? (\textbf{\te{గాయపరచు}}, \emph{gāyaparacu}.) To grapple? (\textbf{\te{పెనుగులాడు}}, \emph{penugulāḍu}.) To gather round? (\textbf{\te{గుమిగూడు}}, \emph{gumigūḍu}.) To throng? (\textbf{\te{ఎగబడు}}, \emph{egabaḍu}.) To fall down? (\textbf{\te{పడిపోవు}}, \emph{paḍipōvu}.) A victim? (\textbf{\te{బాధితుడు}}, \emph{bādhituḍu}.) To petition? (\textbf{\te{మొరపెట్టుకో}}, \emph{morapeṭṭukō}.) To beg? (\textbf{\te{వేడుకో}}, \emph{vēḍukō}.) Ago? (\textbf{\te{క్రితం}}, \emph{kritaṁ}.) Before that? (\textbf{\te{అంతకుముందు}}, \emph{antakumundu}.) From then on? (\textbf{\te{అప్పటినుంచి}}, \emph{appaṭinuñci}.) A dead body? (\textbf{\te{మృతదేహం}}, \emph{mṛtadēhaṁ}.) To haunt? (\textbf{\te{వెంటాడు}}, \emph{veṇṭāḍu}.) To dare? (\textbf{\te{తెగించు}}, \emph{tegiñcu}.) To abduct? (\textbf{\te{అపహరించు}}, \emph{apahariñcu}.) To interrogate? (\textbf{\te{విచారించు}}, \emph{vicāriñcu}.) To confiscate? (\textbf{\te{జప్తు} \te{చేయు}}, \emph{japtu cēyu}.) An invasion? (\textbf{\te{దండయాత్ర}}, \emph{daṇḍayātra}.)
\end{tcolorbox}
\section[(dialogue)]{Daring Deeds — once more}

\label{lesson:TE-R381-review-final-381}



\begin{quote}
Say the Telugu for to dare and to abduct.
\end{quote}

\subsection*{Your turn}
\begin{itemize}
\raggedright
  \item \emph{Say it:} to dare, to abduct, to interrogate, to confiscate, an invasion
\end{itemize}

\begin{tcolorbox}[breakable,colback=teal!4,colframe=teal!35!black,title={Before you move on}]
To dare? (\textbf{\te{తెగించు}}, \emph{tegiñcu}.) To abduct? (\textbf{\te{అపహరించు}}, \emph{apahariñcu}.) To interrogate? (\textbf{\te{విచారించు}}, \emph{vicāriñcu}.) To confiscate? (\textbf{\te{జప్తు} \te{చేయు}}, \emph{japtu cēyu}.) An invasion? (\textbf{\te{దండయాత్ర}}, \emph{daṇḍayātra}.)
\end{tcolorbox}
